% TOP_PROBLEM: {"id": 30006723, "title": "The Mumford-Tate Conjecture for Abelian Varieties over Number Fields", "category_id": 1, "category": {"id": 1, "name": "number_theory", "display_name": "Number Theory", "description": "Properties of integers, prime numbers, Diophantine equations.", "slug": "number-theory", "order_index": 1, "created_at": "2026-07-31T15:26:25.670Z"}, "status": "open"}
% FIRST_POSTED: 2026-09-25
% ENTRY_KIND: statement_only
% !TeX program = lualatex
% Definition and sources only; this file is not an LLM solution attempt.
\documentclass[11pt]{article}
\usepackage[margin=25mm]{geometry}
\usepackage{fontspec}
\setmainfont{Latin Modern Roman}
\usepackage{amsmath,amssymb,mathtools}
\usepackage{unicode-math}
\setmathfont{Latin Modern Math}
\usepackage{xurl,hyperref}
\hypersetup{colorlinks=true,urlcolor=blue}
\setcounter{secnumdepth}{0}
\setlength{\parindent}{0pt}
\setlength{\parskip}{5pt}
\setlength{\emergencystretch}{3em}
\begin{document}
\section*{\texorpdfstring{The Mumford-Tate Conjecture for Abelian Varieties over Number Fields}{The Mumford-Tate Conjecture for Abelian Varieties over Number Fields}}\label{30006723}
\subsection{Definitions and mathematical statement}
For an abelian variety $A/K$, let $V_\ell(A)=T_\ell(A)\otimes{\mathbb{Q}}_\ell$ and let $G_\ell$ be the Zariski closure of the Galois image in $\operatorname{GL}(V_\ell(A))$. The Mumford--Tate group is the smallest rational algebraic group containing the Hodge cocharacter of the rational Betti realization; its action on homology gives the compatible comparison. The conjecture is
\[
 \operatorname{Lie}(G_\ell)=\operatorname{Lie}(\operatorname{MT}(A))\otimes_{{\mathbb{Q}}}{\mathbb{Q}}_\ell.
\]
Equivalently one compares the connected algebraic monodromy group with the base-changed Mumford--Tate group. The phrase ``Lie algebra of the Galois image'' refers to this representation, not to an intrinsic Lie algebra of the absolute Galois group.
\par\smallskip\textit{Formulation and concept references:} \href{https://doi.org/10.1201/9781439864227}{[S1]}.

\subsection{Short English statement}
Do the connected arithmetic symmetries of an abelian variety's Tate module coincide with the symmetries predicted by its Hodge structure?
\subsection{Sources}
\begin{itemize}
\item[S1]J.-P. Serre, Abelian l-Adic Representations and Elliptic Curves, Addison-Wesley, 1968.
\url{https://doi.org/10.1201/9781439864227}.
\textit{Formulation source.}
\item[S2]On the Frobenius fields of abelian varieties over number fields.
\url{https://arxiv.org/html/2402.07935v2}.
\textit{Further reference; not a proof endorsement.}
\end{itemize}
\end{document}
